% Generado por src/comparativa.py — no editar a mano, se regenera.
\begin{table}[htbp]
    \centering
    {\footnotesize
    \begin{tabular}{@{}lrrrrr@{}}
        \toprule
        \textbf{Modelo} & \textbf{media de $Z$} & \textbf{media de $X$} & \textbf{media de $Y$} & \textbf{paridad} & \textbf{corr. de vecinos} \\
        \midrule
        No mitigar & $-$0,055 & $-$0,567 & --- & $-$1,164 & $-$0,282 \\
        f constante & $-$7,751 & $-$11,903 & --- & $-$12,513 & $-$22,667 \\
        Ridge · 1 variable & $-$0,058 & $-$1,527 & --- & $-$0,063 & $-$1,422 \\
        Ridge · 34 variables & $-$3,324 & $-$0,459 & --- & 0,438 & $-$4,957 \\
        Random Forest & 0,004 & $-$0,034 & --- & 0,768 & $-$0,280 \\
        GNN · SAGE & $-$0,048 & 0,226 & --- & $-$0,493 & $-$0,043 \\
        \bottomrule
    \end{tabular}}

    \caption{$R^2$ sobre $f$ restringido a los seis tipos de circuito zero-shot, los que no aparecen en entrenamiento. \texttt{mean\_Y} solo lo define \texttt{random}, así que no tiene columna aquí.}
    \label{tab:comparativa_zeroshot}

    {\footnotesize Fuente: elaboración propia.}
\end{table}
